\documentclass[a4paper,12pt,twoside]{memoir}

% Castellano
\usepackage[spanish,es-tabla]{babel}
\selectlanguage{spanish}
\usepackage[utf8]{inputenc}
\usepackage[T1]{fontenc}
\usepackage{lmodern} % scalable font
\usepackage{microtype}
\usepackage{placeins}

\RequirePackage{booktabs}
\RequirePackage[table]{xcolor}
\RequirePackage{xtab}
\RequirePackage{multirow}

% Links
\PassOptionsToPackage{hyphens}{url}\usepackage[colorlinks]{hyperref}
\hypersetup{
	allcolors = {red}
}

% Ecuaciones
\usepackage{amsmath}

% Rutas de fichero / paquete
\newcommand{\ruta}[1]{{\sffamily #1}}

% Párrafos
\nonzeroparskip

% Huérfanas y viudas
\widowpenalty100000
\clubpenalty100000

% Evitar solapes en el header
\nouppercaseheads

% Imagenes
\usepackage{graphicx}
\newcommand{\imagen}[2]{
	\begin{figure}[!h]
		\centering
		\includegraphics[width=0.9\textwidth]{#1}
		\caption{#2}\label{fig:#1}
	\end{figure}
	\FloatBarrier
}

\newcommand{\imagenflotante}[2]{
	\begin{figure}%[!h]
		\centering
		\includegraphics[width=0.9\textwidth]{#1}
		\caption{#2}\label{fig:#1}
	\end{figure}
}



% El comando \figura nos permite insertar figuras comodamente, y utilizando
% siempre el mismo formato. Los parametros son:
% 1 -> Porcentaje del ancho de página que ocupará la figura (de 0 a 1)
% 2 --> Fichero de la imagen
% 3 --> Texto a pie de imagen
% 4 --> Etiqueta (label) para referencias
% 5 --> Opciones que queramos pasarle al \includegraphics
% 6 --> Opciones de posicionamiento a pasarle a \begin{figure}
\newcommand{\figuraConPosicion}[6]{%
  \setlength{\anchoFloat}{#1\textwidth}%
  \addtolength{\anchoFloat}{-4\fboxsep}%
  \setlength{\anchoFigura}{\anchoFloat}%
  \begin{figure}[#6]
    \begin{center}%
      \Ovalbox{%
        \begin{minipage}{\anchoFloat}%
          \begin{center}%
            \includegraphics[width=\anchoFigura,#5]{#2}%
            \caption{#3}%
            \label{#4}%
          \end{center}%
        \end{minipage}
      }%
    \end{center}%
  \end{figure}%
}

%
% Comando para incluir imágenes en formato apaisado (sin marco).
\newcommand{\figuraApaisadaSinMarco}[5]{%
  \begin{figure}%
    \begin{center}%
    \includegraphics[angle=90,height=#1\textheight,#5]{#2}%
    \caption{#3}%
    \label{#4}%
    \end{center}%
  \end{figure}%
}
% Para las tablas
\newcommand{\otoprule}{\midrule [\heavyrulewidth]}
%
% Nuevo comando para tablas pequeñas (menos de una página).
\newcommand{\tablaSmall}[5]{%
 \begin{table}
  \begin{center}
   \rowcolors {2}{gray!35}{}
   \begin{tabular}{#2}
    \toprule
    #4
    \otoprule
    #5
    \bottomrule
   \end{tabular}
   \caption{#1}
   \label{tabla:#3}
  \end{center}
 \end{table}
}

%
%Para el float H de tablaSmallSinColores
\usepackage{float}

%
% Nuevo comando para tablas pequeñas (menos de una página).
\newcommand{\tablaSmallSinColores}[5]{%
 \begin{table}[H]
  \begin{center}
   \begin{tabular}{#2}
    \toprule
    #4
    \otoprule
    #5
    \bottomrule
   \end{tabular}
   \caption{#1}
   \label{tabla:#3}
  \end{center}
 \end{table}
}

\newcommand{\tablaApaisadaSmall}[5]{%
\begin{landscape}
  \begin{table}
   \begin{center}
    \rowcolors {2}{gray!35}{}
    \begin{tabular}{#2}
     \toprule
     #4
     \otoprule
     #5
     \bottomrule
    \end{tabular}
    \caption{#1}
    \label{tabla:#3}
   \end{center}
  \end{table}
\end{landscape}
}

%
% Nuevo comando para tablas grandes con cabecera y filas alternas coloreadas en gris.
\newcommand{\tabla}[6]{%
  \begin{center}
    \tablefirsthead{
      \toprule
      #5
      \otoprule
    }
    \tablehead{
      \multicolumn{#3}{l}{\small\sl continúa desde la página anterior}\\
      \toprule
      #5
      \otoprule
    }
    \tabletail{
      \hline
      \multicolumn{#3}{r}{\small\sl continúa en la página siguiente}\\
    }
    \tablelasttail{
      \hline
    }
    \bottomcaption{#1}
    \rowcolors {2}{gray!35}{}
    \begin{xtabular}{#2}
      #6
      \bottomrule
    \end{xtabular}
    \label{tabla:#4}
  \end{center}
}

%
% Nuevo comando para tablas grandes con cabecera.
\newcommand{\tablaSinColores}[6]{%
  \begin{center}
    \tablefirsthead{
      \toprule
      #5
      \otoprule
    }
    \tablehead{
      \multicolumn{#3}{l}{\small\sl continúa desde la página anterior}\\
      \toprule
      #5
      \otoprule
    }
    \tabletail{
      \hline
      \multicolumn{#3}{r}{\small\sl continúa en la página siguiente}\\
    }
    \tablelasttail{
      \hline
    }
    \bottomcaption{#1}
    \begin{xtabular}{#2}
      #6
      \bottomrule
    \end{xtabular}
    \label{tabla:#4}
  \end{center}
}

%
% Nuevo comando para tablas grandes sin cabecera.
\newcommand{\tablaSinCabecera}[5]{%
  \begin{center}
    \tablefirsthead{
      \toprule
    }
    \tablehead{
      \multicolumn{#3}{l}{\small\sl continúa desde la página anterior}\\
      \hline
    }
    \tabletail{
      \hline
      \multicolumn{#3}{r}{\small\sl continúa en la página siguiente}\\
    }
    \tablelasttail{
      \hline
    }
    \bottomcaption{#1}
  \begin{xtabular}{#2}
    #5
   \bottomrule
  \end{xtabular}
  \label{tabla:#4}
  \end{center}
}



\definecolor{cgoLight}{HTML}{EEEEEE}
\definecolor{cgoExtralight}{HTML}{FFFFFF}

%
% Nuevo comando para tablas grandes sin cabecera.
\newcommand{\tablaSinCabeceraConBandas}[5]{%
  \begin{center}
    \tablefirsthead{
      \toprule
    }
    \tablehead{
      \multicolumn{#3}{l}{\small\sl continúa desde la página anterior}\\
      \hline
    }
    \tabletail{
      \hline
      \multicolumn{#3}{r}{\small\sl continúa en la página siguiente}\\
    }
    \tablelasttail{
      \hline
    }
    \bottomcaption{#1}
    \rowcolors[]{1}{cgoExtralight}{cgoLight}

  \begin{xtabular}{#2}
    #5
   \bottomrule
  \end{xtabular}
  \label{tabla:#4}
  \end{center}
}




\graphicspath{ {./img/} }

% Capítulos
\chapterstyle{bianchi}
\newcommand{\capitulo}[2]{
	\setcounter{chapter}{#1}
	\setcounter{section}{0}
	\setcounter{figure}{0}
	\setcounter{table}{0}
	\chapter*{#2}
	\addcontentsline{toc}{chapter}{#2}
	\markboth{#2}{#2}
}

% Apéndices
\renewcommand{\appendixname}{Apéndice}
\renewcommand*\cftappendixname{\appendixname}

\newcommand{\apendice}[1]{
	%\renewcommand{\thechapter}{A}
	\chapter{#1}
}

\renewcommand*\cftappendixname{\appendixname\ }

% Formato de portada
\makeatletter
\usepackage{xcolor}
\newcommand{\tutor}[1]{\def\@tutor{#1}}
\newcommand{\course}[1]{\def\@course{#1}}
\definecolor{cpardoBox}{HTML}{E6E6FF}
\def\maketitle{
  \null
  \thispagestyle{empty}
  % Cabecera ----------------
\noindent\includegraphics[width=\textwidth]{cabecera}\vspace{1cm}%
  \vfill
  % Título proyecto y escudo informática ----------------
  \colorbox{cpardoBox}{%
    \begin{minipage}{.8\textwidth}
      \vspace{.5cm}\Large
      \begin{center}
      \textbf{TFG del Grado en Ingeniería Informática}\vspace{.6cm}\\
      \textbf{\LARGE\@title{}}
      \end{center}
      \vspace{.2cm}
    \end{minipage}

  }%
  \hfill\begin{minipage}{.20\textwidth}
    \includegraphics[width=\textwidth]{escudoInfor}
  \end{minipage}
  \vfill
  % Datos de alumno, curso y tutores ------------------
  \begin{center}%
  {%
    \noindent\LARGE
    Presentado por \@author{}\\ 
    en Universidad de Burgos --- \@date{}\\
    Tutor: \@tutor{}\\
  }%
  \end{center}%
  \null
  \cleardoublepage
  }
\makeatother


% Datos de portada
\title{Desarrollo y explotación de
	un chatbot de FAQs sobre
	una asignatura de Trabajo
	Fin de Grado \\UBU-Chatbot}
\author{Alfredo Asensio Vázquez}
\tutor{Dr. Raúl Marticorena Sánchez}
\date{julio de 2021}

\begin{document}

\maketitle



\cleardoublepage



%%%%%%%%%%%%%%%%%%%%%%%%%%%%%%%%%%%%%%%%%%%%%%%%%%%%%%%%%%%%%%%%%%%%%%%%%%%%%%%%%%%%%%%%



\frontmatter


\clearpage

% Indices
\tableofcontents

\clearpage

\listoffigures

\clearpage

\listoftables

\clearpage

\mainmatter

\appendix

\apendice{Plan de Proyecto Software}

\section{Introducción}

\section{Planificación temporal}

\section{Estudio de viabilidad}

\subsection{Viabilidad económica}

\subsection{Viabilidad legal}



\apendice{Especificación de Requisitos}

\section{Introducción}

En este apartado se detallan los objetivos generales del proyecto, los requisitos y la especificación de los mismos.  
Se analizan tanto los requisitos funcionales como los no funcionales. 
\begin{itemize}
	\tightlist
	\item
	Requisitos funcionales: comportamientos específicos que debe tener el sistema. Se relacionan con los casos de uso.
	\item
	Requisitos no funcionales: especifican los criterios a seguir, restricciones y condiciones que impone el cliente.

\end{itemize}

\section{Objetivos generales}


\begin{itemize}
	\tightlist
	\item
	Implementar un \textit{chatbot} que responda las preguntas frecuentes de la asignatura Trabajo de Fin de Grado tanto para la modalidad online como presencial.
	\item
	Ser capaces de dar respuesta al mayor número de preguntas posible.
	\item
	Integrar el \textit{chatbot} en UBUVirtual y Slack.
	\item
	Diseñar una interfaz conversacional que resulte agradable a los usuarios y respete la imagen corporativa.
	\item
	Capturar los \textit{logs} de las conversaciones para analizarlos y extraer resultados. 
\end{itemize}
\newpage

\section{Catalogo de requisitos}


\subsection{Requisitos funcionales}\label{requisitosFuncionales}

\begin{itemize}
	\tightlist
	\item
	\textbf{RF-1 Interacción texual:} la aplicación debe permitir al usuario interactuar con ella mediante texto.
	\item
	\textbf{RF-2 Reconocimiento de preguntas:} la aplicación debe ser capaz de reconocer las preguntas que introduce el usuario.
	\item
	\textbf{RF-3 Respuesta a las preguntas:} la aplicación debe ser capaz de responder a las preguntas que le formule el usuario.
	\item
	\textbf{RF-4 Información mediante hipervínculos:} la aplicación debe adjuntar como hipervínculos las redirecciones a otras páginas.
	\item
	\textbf{RF-5 Minimizar la interfaz conversacional:} la aplicación debe permitir minimizar la interfaz conversacional.
	\item
	\textbf{RF-6 Expandir la interfaz conversacional:} la aplicación debe permitir expandir la interfaz conversacional.
	\item
	\textbf{RF-7 Informe de error:} la aplicación debe informar de un error al usuario cuando su mensaje de entrada no sea un texto válido o no haya sido posible interpretarlo. 
\end{itemize}


\subsection{Requisitos no funcionales}\label{requisitosNoFuncionales}

\begin{itemize}
	\tightlist
	\item
	\textbf{RNF-1 Rendimiento:} la aplicación tiene que tener un tiempo de respuesta bajo.  
	\item
	\textbf{RNF-2 Usabilidad:} la aplicación debe ser intuitiva y fácil de entender y utilizar.  
	\item
	\textbf{RNF-3 Colores corporativos:} la aplicación debe mantener los colores corporativos de UBUVirtual.  
	\item
	\textbf{RNF-4 Disponibilidad:} la aplicación debe estar disponible el mayor tiempo posible.
	\item
	\textbf{RNF-5 Mantenibilidad:} la aplicación debe ser fácilmente modificable.
	\item
	\textbf{RNF-6 Registro de actividad:} la aplicación debe almacenar las interacciones llevadas a cabo con los usuarios.
	\item
	\textbf{RNF-7 Portabilidad:} la aplicación debe almacenar las interacciones llevadas a cabo con los usuarios.
	\item
	\textbf{RNF-8 Compatibilidad de navegadores:} la aplicación debe ser compatible para los navegadores más importantes, incluyendo también los navegadores de \textit{smartphones}.
\end{itemize}

\newpage
\section{Especificación de requisitos}

\subsection{Diagrama de casos de uso}
\imagen{diagramaCasosDeUso}{Diagrama de casos de uso.}

El actor -usuario- va a ser el estudiante, profesor o cualquier persona con acceso a la asignatura que comience una nueva sesión con el \textit{chatbot}.

\newpage

\subsection{Casos de uso}

\begin{longtable}[H]{@{}ll@{}}
	\toprule
	\begin{minipage}[b]{0.23\columnwidth}\raggedright\strut
		\textbf{CU-01}\strut
	\end{minipage} & \begin{minipage}[b]{0.71\columnwidth}\raggedright\strut
		\textbf{Expandir la interfaz conversacional.}\strut
	\end{minipage}\tabularnewline
	\midrule
	\endhead  
	\begin{minipage}[t]{0.23\columnwidth}\raggedright\strut
		\textbf{Requisitos asociados}\strut
	\end{minipage} & \begin{minipage}[t]{0.71\columnwidth}\raggedright\strut
		RF-6\strut
	\end{minipage}\tabularnewline
	\begin{minipage}[t]{0.23\columnwidth}\raggedright\strut
		\textbf{Descripción}\strut
	\end{minipage} & \begin{minipage}[t]{0.71\columnwidth}\raggedright\strut
		Permite al usuario expandir la interfaz conversacional del \textit{chatbot}.\strut
	\end{minipage}\tabularnewline
	\begin{minipage}[t]{0.23\columnwidth}\raggedright\strut
		\textbf{Precondición}\strut
	\end{minipage} & \begin{minipage}[t]{0.71\columnwidth}\raggedright\strut
		El usuario está logeado en UBUVirtual y dentro de la asignatura Trabajo Fin de Grado, en el apartado \textit{Chatbot de preguntas frecuentes}.\\
		El \textit{chatbot} esta minimizado.\strut
	\end{minipage}\tabularnewline
	\begin{minipage}[t]{0.23\columnwidth}\raggedright\strut
		\textbf{Acciones}\strut
	\end{minipage} & \begin{minipage}[t]{0.71\columnwidth}\raggedright\strut
		\begin{enumerate}
			\def\labelenumi{\arabic{enumi}.}
			\tightlist
			\item
			El usuario accede a la página en que aparece el icono del \textit{chatbot}.
			\item
			El usuario hace clic en el botón circular con el logo del proyecto.
			\item
			Se expande la interfaz conversacional.
		\end{enumerate}\strut
	\end{minipage}\tabularnewline
	\begin{minipage}[t]{0.23\columnwidth}\raggedright\strut
		\textbf{Postcondición}\strut
	\end{minipage} & \begin{minipage}[t]{0.71\columnwidth}\raggedright\strut
		La interfaz conversacional queda expandida y el usuario puede interactuar con el \textit{chatbot}.\strut
	\end{minipage}\tabularnewline
	\begin{minipage}[t]{0.23\columnwidth}\raggedright\strut
		\textbf{Excepciones}\strut
	\end{minipage} & \begin{minipage}[t]{0.71\columnwidth}\raggedright\strut
		-\strut
	\end{minipage}\tabularnewline
	\begin{minipage}[t]{0.23\columnwidth}\raggedright\strut
		\textbf{Importancia}\strut
	\end{minipage} & \begin{minipage}[t]{0.71\columnwidth}\raggedright\strut
		Alta\strut
	\end{minipage}\tabularnewline
	\bottomrule
	\caption{CU-01 Expandir la interfaz conversacional.}
\end{longtable}

\newpage
\begin{longtable}[H]{@{}ll@{}}
	\toprule
	\begin{minipage}[b]{0.23\columnwidth}\raggedright\strut
		\textbf{CU-02}\strut
	\end{minipage} & \begin{minipage}[b]{0.71\columnwidth}\raggedright\strut
		\textbf{Formular una pregunta.}\strut
	\end{minipage}\tabularnewline
	\midrule
	\endhead  
	\begin{minipage}[t]{0.23\columnwidth}\raggedright\strut
		\textbf{Requisitos asociados}\strut
	\end{minipage} & \begin{minipage}[t]{0.71\columnwidth}\raggedright\strut
		RF-1, RF-2, RF-3, RF-4, RF-7\strut
	\end{minipage}\tabularnewline
	\begin{minipage}[t]{0.23\columnwidth}\raggedright\strut
		\textbf{Descripción}\strut
	\end{minipage} & \begin{minipage}[t]{0.71\columnwidth}\raggedright\strut
		El usuario introduce de manera textual una pregunta al \textit{chatbot}.\strut
	\end{minipage}\tabularnewline
	\begin{minipage}[t]{0.23\columnwidth}\raggedright\strut
		\textbf{Precondición}\strut
	\end{minipage} & \begin{minipage}[t]{0.71\columnwidth}\raggedright\strut
		El usuario está logeado en UBUVirtual y dentro de la asignatura Trabajo Fin de Grado, en el apartado \textit{Chatbot de preguntas frecuentes}.\\
		El \textit{chatbot} esta expandido.\strut
	\end{minipage}\tabularnewline
	\begin{minipage}[t]{0.23\columnwidth}\raggedright\strut
		\textbf{Acciones}\strut
	\end{minipage} & \begin{minipage}[t]{0.71\columnwidth}\raggedright\strut
		\begin{enumerate}
			\def\labelenumi{\arabic{enumi}.}
			\tightlist
			\item
			El usuario escribe y envía por medio de la interfaz conversacional una pregunta al \textit{chatbot}.
			\item
			El programa analiza la pregunta introducida e intenta reconocerla.
			\item
			Se responde al usuario en modo texto y por medio de la interfaz conversacional a su pregunta. 
			\item
			El \textit{chatbot} queda a la espera de recibir nuevas preguntas.
		\end{enumerate}\strut
	\end{minipage}\tabularnewline
	\begin{minipage}[t]{0.23\columnwidth}\raggedright\strut
		\textbf{Postcondición}\strut
	\end{minipage} & \begin{minipage}[t]{0.71\columnwidth}\raggedright\strut
		Se devuelve un mensaje con la respuesta a la pregunta.\strut
	\end{minipage}\tabularnewline
	\begin{minipage}[t]{0.23\columnwidth}\raggedright\strut
		\textbf{Excepciones}\strut
	\end{minipage} & \begin{minipage}[t]{0.71\columnwidth}\raggedright\strut
		Si el \textit{chatbot} no es capaz de reconocer la pregunta introducida por el usuario o está mal formulada se informa al usuario por medio de un mensaje de que no ha sido posible entenderle.\strut
	\end{minipage}\tabularnewline
	\begin{minipage}[t]{0.23\columnwidth}\raggedright\strut
		\textbf{Importancia}\strut
	\end{minipage} & \begin{minipage}[t]{0.71\columnwidth}\raggedright\strut
		Alta\strut
	\end{minipage}\tabularnewline
	\bottomrule
	\caption{CU-02 Formular una pregunta.}
\end{longtable}

\newpage
\begin{longtable}[H]{@{}ll@{}}
	\toprule
	\begin{minipage}[b]{0.23\columnwidth}\raggedright\strut
		\textbf{CU-03}\strut
	\end{minipage} & \begin{minipage}[b]{0.71\columnwidth}\raggedright\strut
		\textbf{Formular una pregunta. (Slack) }\strut
	\end{minipage}\tabularnewline
	\midrule
	\endhead  
	\begin{minipage}[t]{0.23\columnwidth}\raggedright\strut
		\textbf{Requisitos asociados}\strut
	\end{minipage} & \begin{minipage}[t]{0.71\columnwidth}\raggedright\strut
		RF-1, RF-2, RF-3, RF-4, RF-7\strut
	\end{minipage}\tabularnewline
	\begin{minipage}[t]{0.23\columnwidth}\raggedright\strut
		\textbf{Descripción}\strut
	\end{minipage} & \begin{minipage}[t]{0.71\columnwidth}\raggedright\strut
		El usuario introduce de manera textual una pregunta al \textit{chatbot}.\strut
	\end{minipage}\tabularnewline
	\begin{minipage}[t]{0.23\columnwidth}\raggedright\strut
		\textbf{Precondición}\strut
	\end{minipage} & \begin{minipage}[t]{0.71\columnwidth}\raggedright\strut
		El usuario está logeado en el espacio de trabajo de Slack y en la aplicación \textit{UBU Asistante Virtual}.
	\end{minipage}\tabularnewline
	\begin{minipage}[t]{0.23\columnwidth}\raggedright\strut
		\textbf{Acciones}\strut
	\end{minipage} & \begin{minipage}[t]{0.71\columnwidth}\raggedright\strut
		\begin{enumerate}
			\def\labelenumi{\arabic{enumi}.}
			\tightlist
			\item
			El usuario escribe y envía por medio de la interfaz conversacional una pregunta al \textit{chatbot	}.
			\item
			El programa analiza la pregunta introducida e intenta reconocerla.
			\item
			Se responde al usuario en modo texto y por medio de la interfaz conversacional a su pregunta. 
			\item
			El \textit{chatbot} queda a la espera de recibir nuevas preguntas.
		\end{enumerate}\strut
	\end{minipage}\tabularnewline
	\begin{minipage}[t]{0.23\columnwidth}\raggedright\strut
		\textbf{Postcondición}\strut
	\end{minipage} & \begin{minipage}[t]{0.71\columnwidth}\raggedright\strut
		Se devuelve un mensaje con la respuesta a la pregunta.\strut
	\end{minipage}\tabularnewline
	\begin{minipage}[t]{0.23\columnwidth}\raggedright\strut
		\textbf{Excepciones}\strut
	\end{minipage} & \begin{minipage}[t]{0.71\columnwidth}\raggedright\strut
		Si el \textit{chatbot} no es capaz de reconocer la pregunta introducida por el usuario o está mal formulada se informa al usuario por medio de un mensaje de que no ha sido posible entenderle.\strut
	\end{minipage}\tabularnewline
	\begin{minipage}[t]{0.23\columnwidth}\raggedright\strut
		\textbf{Importancia}\strut
	\end{minipage} & \begin{minipage}[t]{0.71\columnwidth}\raggedright\strut
		Alta\strut
	\end{minipage}\tabularnewline
	\bottomrule
	\caption{CU-03 Formular una pregunta. (Slack)}
\end{longtable}

\newpage


\begin{longtable}[H]{@{}ll@{}}
	\toprule
	\begin{minipage}[b]{0.23\columnwidth}\raggedright\strut
		\textbf{CU-04}\strut
	\end{minipage} & \begin{minipage}[b]{0.71\columnwidth}\raggedright\strut
		\textbf{Ocultar interfaz conversacional.}\strut
	\end{minipage}\tabularnewline
	\midrule
	\endhead  
	\begin{minipage}[t]{0.23\columnwidth}\raggedright\strut
		\textbf{Requisitos asociados}\strut
	\end{minipage} & \begin{minipage}[t]{0.71\columnwidth}\raggedright\strut
		RF-5\strut
	\end{minipage}\tabularnewline
	\begin{minipage}[t]{0.23\columnwidth}\raggedright\strut
		\textbf{Descripción}\strut
	\end{minipage} & \begin{minipage}[t]{0.71\columnwidth}\raggedright\strut
		Permite al usuario ocultar la interfaz conversacional del \textit{chatbot}.\strut
	\end{minipage}\tabularnewline
	\begin{minipage}[t]{0.23\columnwidth}\raggedright\strut
		\textbf{Precondición}\strut
	\end{minipage} & \begin{minipage}[t]{0.71\columnwidth}\raggedright\strut
		El usuario está logeado en UBUVirtual y dentro de la asignatura Trabajo Fin de Grado, en el apartado \textit{Chatbot de preguntas frecuentes}.\\
		El \textit{chatbot} esta expandido.\strut
	\end{minipage}\tabularnewline
	\begin{minipage}[t]{0.23\columnwidth}\raggedright\strut
		\textbf{Acciones}\strut
	\end{minipage} & \begin{minipage}[t]{0.71\columnwidth}\raggedright\strut
		\begin{enumerate}
			\def\labelenumi{\arabic{enumi}.}
			\tightlist
			\item 
			El usuario hace clic en el botón circular con el logo del proyecto.
			\item
			La interfaz conversacional queda oculta.
		\end{enumerate}\strut
	\end{minipage}\tabularnewline
	\begin{minipage}[t]{0.23\columnwidth}\raggedright\strut
		\textbf{Postcondición}\strut
	\end{minipage} & \begin{minipage}[t]{0.71\columnwidth}\raggedright\strut
		La interfaz conversacional queda oculta.\strut
	\end{minipage}\tabularnewline
	\begin{minipage}[t]{0.23\columnwidth}\raggedright\strut
		\textbf{Excepciones}\strut
	\end{minipage} & \begin{minipage}[t]{0.71\columnwidth}\raggedright\strut
		-\strut
	\end{minipage}\tabularnewline
	\begin{minipage}[t]{0.23\columnwidth}\raggedright\strut
		\textbf{Importancia}\strut
	\end{minipage} & \begin{minipage}[t]{0.71\columnwidth}\raggedright\strut
		Alta\strut
	\end{minipage}\tabularnewline
	\bottomrule
	\caption{CU-04 Ocultar interfaz conversacional.}
\end{longtable}
\apendice{Especificación de diseño}

\section{Introducción}

En este apartado se explica con detalle como se han estructurado los datos, los procedimientos y la arquitectura que se ha utilizado.

\section{Diseño de datos}

Los datos se estructuran en torno a los \textit{intents} y \textit{entities}.

\subsection{Entities}

Hemos considerado las tres entidades mostradas en la figura \ref{fig:Entities}, al darse el caso de que había sinónimos de estos que se repetían en muchas ocasiones.

\imagenEnana{Entities}{Entidades utilizadas en el agente.}

\subsection{Intents}

Aquí es donde almacenamos nuestras frases de entrenamiento y respuestas. Disponemos de un total de 205 archivos JSON para 103 intents diferentes en la modalidad online, y 195 archivos para 98 \textit{intents} diferentes en la modalidad presencial. En las Figuras \ref{fig:bbdd1} y \ref{fig:bbdd2} se presenta el esquema utilizado para la versión online.

\imagenGrande{bbdd1}{Diseño de \textit{intents} del agente para la modalidad online (Parte 1).}

\imagenGrande{bbdd2}{Diseño de \textit{intents} del agente para la modalidad online (Parte 2).}

Se ha utilizado la nomenclatura de comenzar el nombre del \textit{intent} con un simbolo \textit{\_} para determinar aquellos cuya respuesta es una \textit{rich response}.

\newpage
Para hacer legible y mantenible nuestra base de datos de preguntas y respuestas se dispone de un documento de texto en el que se recopila toda la información. De esta manera una persona ajena al proyecto podrá comprender la estructura rápidamente.

\imagenGrande{recopilacion}{Primera página del documento recopilatorio de preguntas y respuestas.}


\section{Diseño procedimental}

\section{Diseño arquitectónico}



\apendice{Documentación técnica de programación}

\section{Introducción}

\section{Estructura de directorios}

\section{Manual del programador}

\section{Compilación, instalación y ejecución del proyecto}

\section{Pruebas del sistema}

\apendice{Documentación de usuario}

\section{Introducción}

En este apartado se explican los requerimientos de la aplicación para ser ejecutada tanto en la plataforma Moodle como en Slack. Se indica también el proceso de instalación en ambas plataformas. Finalmente, en el manual de usuario se dan las indicaciones para utilizar correctamente la aplicación.

\section{Requisitos de usuarios}

En esta sección se indican los requisitos para utilizar el chatbot en cada plataforma para la versión de la modalidad online.

\subsection{Plataforma UBU Virtual}

\begin{itemize}
	\tightlist
	\item
	Conexión a Internet.
	\item
	Navegador web en cualquier dispositivo: PC, Smartphone, Tablet, etc. \\
	Compatible con Chrome, Safari, Firefox, Edge y Opera entre otros. Incompatible con Internet Explorer.
	\item
	Cuenta activa en la plataforma UBUVirtual.
	\item
	Permiso de acceso al curso `TRABAJO FIN DE GRADO (Grupo 90) On-line' (id 11707).
	
\end{itemize}

\newpage
\subsection{Slack}	

\begin{itemize}
	\tightlist
	\item
	Conexión a Internet.
	\item
	Navegador web en cualquier dispositivo: PC, Smartphone, Tablet, etc. \\
	Compatible con Chrome, Safari, Firefox, Edge y Opera. \cite{SlackRequisitos}
	\item
	Cuenta de Slack y permiso de acceso al espacio de trabajo ubuchatbot.slack.com. 

\end{itemize}
 
 \newpage
\section{Instalación}

\subsection{Plataforma UBUVirtual}

No es necesario ningún proceso de instalación. Únicamente necesitaremos acceder a la asignatura en la plataforma UBUVirtual con un navegador compatible y tener JavaScript activado.\\

El enlace de acceso al chatbot es el siguiente: \url{https://ubuvirtual.ubu.es/mod/page/view.php?id=3132682}

\subsection{Slack}

Se puede acceder a esta plataforma mediante la versión web o la aplicación de escritorio.\\
En caso de que quisiéramos obtener la versión de escritorio, deberíamos acceder a \url{https://slack.com/intl/es-es/downloads} y seleccionar la versión compatible con nuestro sistema operativo, pero este paso no es necesario.

Como paso previo necesitaremos una invitación al espacio de trabajo que nos deberá haber enviado alguno de los miembros. En caso de que no se te haya facilitado esta invitación, escribe un email al administrador a su correo  aav0038@alu.ubu.es indicando tu correo de Google y te la hará llegar. \\

Te llegará un email a la bandeja de entrada con el título `Alfredo te ha invitado a trabajar en Slack'. En caso de no recibirlo se debe comprobar la bandeja de correo no deseado. 

\imagenMediana{Slack-EmailConfirmacion}{Email de confirmación una vez recibido acceso}

Para acceder a la aplicación, lo haremos por medio del botón morado de `Únete ahora' disponible en el email de confirmación como se muestra en la Figura \ref{fig:Slack-EmailConfirmacion}. \\
Nos llevará a una pantalla de login en la que al haber introducido un correo de Google, haremos clic en `Continuar con Google' como se muestra en la Figura \ref{fig:Slack-Unirse}.

\imagenMediana{Slack-Unirse}{Pantalla para unirse al espacio de trabajo de Slack}

Seleccionaremos nuestra cuenta de Google y haremos clic en el botón morado `Crear cuenta' como se muestra en la Figura \ref{fig:Slack-Unirse2}.

\imagenMediana{Slack-Unirse2}{Confirmación para crear la cuenta y acceder al espacio de trabajo.}

Tras unos segundos de carga habremos accedido al espacio de trabajo `UBU Chatbot'.

Alternativamente y en posteriores conexiones, se puede acceder a través del siguiente enlace que nos llevará al espacio de trabajo en el que está integrado el chatbot:\\
\url{https://ubuchatbot.slack.com}
\imagen{anexo-E3-1}{Espacio de trabajo de Slack}

Se nos pedirá credenciales para acceder de la manera que se muestra en la Figura \ref{fig:anexo-E3-1} y lo haremos de la  forma anteriormente explicada.

\newpage

\section{Manual del usuario}

\subsection{Plataforma UBU Virtual}

\subsubsection{Desde PC} \label{subsection:marca-manualpc}

Para la versión de ordenador, una vez hemos accedido a la página de UBUVirtual donde está disponible el chatbot, veremos como en la parte inferior derecha aparece un icono redondo con el logo del proyecto: un oso con un bocadillo de texto sobre un fondo granate. A los segundos aparece un mensaje de bienvenida junto a este, tal y como se muestra en la Figura \ref{fig:anexo-E4-1}.

\imagen{anexo-E4-1}{Visualización del chatbot inicialmente minimizado.}

Para abrir el chatbot haremos clic en el círculo del logo o en el mensaje de bienvenida. Se nos desplegará el chat, en el que aparecerá el mensaje de bienvenida y se habilitará una entrada de texto en la que deberemos hacer clic para empezar a introducir nuestros mensajes. Ver Figura \ref{fig:anexo-E4-2}.

\imagen{anexo-E4-2}{Ventana del chatbot desplegada.}

Podemos comenzar realizando cualquier pregunta. Es importante tener en cuenta que el chatbot esta diseñado para dar respuesta a \textbf{estudiantes online del TFG}. Es decir, si le hacemos preguntas como si fuésemos profesores de la asignatura no las va a saber responder, así como si le hacemos preguntas sobre otras asignaturas. Además, todas las respuestas las va a dar enfocadas a la modalidad online, que pueden no coincidir en muchos casos con la que se daría para presencial.\\
Comenzamos por ejemplo saludando al chatbot, aunque es un paso innecesario.

\imagenPequena{anexo-E4-3}{Respuesta a saludo.}

Como vemos en la Figura \ref{fig:anexo-E4-3} al no haber introducido una pregunta nos devolverá el mismo saludo inicial y nos volverá a pedir que le introduzcamos una pregunta.\\
En este caso vamos a formular la primera pregunta. Podemos hacer cualquier pregunta relativa al Trabajo de Fin de Grado. 

\imagenPequena{anexo-E4-4}{Pregunta bien formulada.}

En este primer ejemplo mostrado en la Figura \ref{fig:anexo-E4-4} se ha formulado la pregunta con signos de interrogación y sin fallos ortográficos. Esto no es un requisito necesario para que el chatbot sea capaz de responder. Podemos formular las preguntas de muchas maneras diferentes.


\imagenPequena{anexo-E4-5}{Distintas formulaciones.}

Como se ve en la Figura \ref{fig:anexo-E4-5}, distintas formulaciones de la misma pregunta han producido la misma respuesta, ya que el chatbot ha sido capaz de comprenderlas. \\
No obstante, hay que seguir una serie de buenas prácticas para que el chatbot funcione de la mejor manera posible:
\begin{itemize}
	\tightlist
	\item
	Solo una pregunta por mensaje.
	\item
	El mensaje ha de ser lo más concreto posible.
	\item
	La pregunta no debe dar lugar a ambigüedades. Ejemplo: si escribimos ``memoria'' no le estamos dando información suficiente para saber qué información queremos conocer acerca de la memoria.
	\item
	No hacer preguntas personales. Si le preguntamos cuál es nuestra hora de defensa no va a saber responderlo. Hay que recordar que no estamos identificados (el chatbot no sabe qué alumno somos) y solo se dan respuestas generales a todos los alumnos.
\end{itemize}


\imagenPequena{anexo-E4-6}{Pregunta mal formulada.}

En la Figura \ref{fig:anexo-E4-6} vemos un ejemplo de pregunta mal formulada. Buscaba el mismo objetivo que las anteriores, pero en este caso el chatbot no ha sido capaz de identificar que información quería conocer el alumno. Esto se debe a que la pregunta no ha sido concreta y ha introducido información que no tenía nada que ver con la pregunta. En este caso la respuesta ha sido un mensaje de error en el que se aportan indicaciones para formular mejor la pregunta y aumentar así las probabilidades de éxito.

\imagenPequena{anexo-E4-7}{Cerrar chatbot.}

La conversación continuará indefinidamente. Cuando deseemos cerrar la conversación simplemente hacemos clic en el círculo con la `x' en la parte inferior del chatbot, señalizado con la flecha roja en la Figura \ref{fig:anexo-E4-7}.


\imagen{anexo-E4-8}{Chatbot cerrado.}

Al hacer clic en el botón de cerrar se volverá al estado inicial en el que la conversación no está desplegada. En este caso, el mensaje de bienvenida ya no aparecerá y únicamente se verá el círculo con el logo del proyecto con el que haciendo clic podremos retomar la conversación cuando queramos. Ver Figura \ref{fig:anexo-E4-8}\\
Si cambiamos de página o refrescamos se reiniciará la conversación, borrándose los mensajes anteriores.
 
\newpage
\subsubsection{Smartphone}

El funcionamiento en Smartphone es muy similar al de PC, primero abrimos la página en la que está integrado el chatbot. 

\imagenPequena{anexo-E4-9}{Página del chatbot desde Safari en iOS.}

En la Figura \ref{fig:anexo-E4-9} podemos ver como en esta versión también aparece un pequeño icono con el logo del proyecto, esta vez sin que se despliegue el mensaje de bienvenida. Hacemos clic en dicho botón.

\newpage

Se nos habrá desplegado una ventana de chat algo diferente a la de versión PC, ya que en este caso ocupa toda la pantalla. Ver Figura \ref{fig:anexo-E4-10}

\imagenPequena{anexo-E4-10}{Ventana del chatbot en Safari iOS.}

Hacemos clic en el cuadro de texto inferior con el texto sombreado ``Haz una pregunta...''.

\newpage

Al clicar nos aparecerá el teclado del teléfono móvil como se muestra en la Figura \ref{fig:anexo-E4-11}, y ya podremos teclear nuestra pregunta. Una vez la tengamos escrita hacemos clic en el botón `intro' del teclado. Este botón puede ser distinto en otros dispositivos móviles.

\imagenPequena{anexo-E4-11}{Página chatbot desde Safari en iOS.}

El funcionamiento y recomendaciones para realizar las preguntas es el mismo que el explicado en la versión de PC, por lo que es recomendable leer dicha sección.

\newpage

Una vez hemos hecho clic en `intro' y enviado nuestro mensaje recibiremos la respuesta del chatbot.
\imagenPequena{anexo-E4-12}{Respuesta del chatbot y botón de cierre de teclado desde Safari en iOS.}

El zoom de la pantalla puede haber sufrido modificaciones, por lo que se recomienda ajustarlo a nuestro gusto.\\
Para cerrar el chat deberemos primero hacer clic en el botón de `ok' marcado en el cuadro rojo en la Figura \ref{fig:anexo-E4-12}, lo cual cerrará el teclado.

\newpage

Una vez se nos ha cerrado el teclado nos aparecerá una fila de opciones en la parte inferior de la pantalla. Hacemos clic en la flecha de volver marcada en rojo en la Figura \ref{fig:anexo-E4-13} para salir definitivamente del chat y volver a la página anterior.

\imagenPequena{anexo-E4-13}{Volver a la página anterior desde Safari en iOS.}

Deslizando la pantalla hacia la izquierda también se podía haber hecho este proceso sin necesidad de cerrar previamente el teclado.\\
Nos devolverá a la página principal de la asignatura.


\newpage
\subsection{Slack}

Accedemos al espacio de trabajo en el que está integrado el chatbot como se ha explicado en el apartado de instalación.\\
URL: \url{https://ubuchatbot.slack.com}

\imagen{anexo-E4-14}{Página principal del espacio de trabajo de Slack.}

Hacemos clic dentro de la sección de `Aplicaciones' en `UBU Asistente Virtual', marcado en la Figura \ref{fig:anexo-E4-14} con un rectángulo rojo.

\newpage

En la página que se nos habrá abierto nos aparecerá el mensaje de bienvenida del chatbot igual que se muestra en la figura \ref{fig:anexo-E4-15}. Ya podemos empezar a introducir nuestras preguntas de la misma forma que en el chatbot de UBUVirtual. Es recomendable seguir las mismas recomendaciones que se han dado en la primera parte del apartado \ref{subsection:marca-manualpc} .

\imagen{anexo-E4-15}{Página del chatbot en Slack.}




\bibliographystyle{plain}
\bibliography{bibliografiaAnexos}

\end{document}
